\documentclass[12pt,a4paper]{article}
\input{../../.common/preamble.sty}

\title{\textbf{La palabra más larga}}

\begin{document}
\maketitle

Dado un fichero de texto encuentra la \textbf{palabra más larga}.

\subsection*{Notas}

\begin{itemize}
    \item Si hay varias palabras con la misma longitud devuelve la última ocurrencia.
    \item Ten en cuenta (únicamente) los siguientes símbolos como frontera de palabras:\\\textcolor{magenta}{\texttt{,.;:()}}
\end{itemize}

\subsection*{Ejemplo}

\textit{Entrada}:
\begin{minted}[
    linenos,
    numbersep=5pt,
    frame=single,
    framesep=3mm,
]{text}
Python es un lenguaje de programación de alto nivel, interpretado y de
tipado dinámico. Fue creado por Guido van Rossum y lanzado por primera
vez en 1991. Su diseño enfatiza la legibilidad del código y la
simplicidad, lo que lo hace accesible tanto para principiantes como para
programadores experimentados. La sintaxis de Python es clara y concisa,
lo que permite escribir código limpio y fácil de mantener.
\end{minted}

\textit{Salida}:\\

La palabra más larga es: \textbf{experimentados}

\end{document}
